% 强弱对偶 + 对偶 gap：横轴对偶变量，纵轴目标值
\documentclass[tikz,border=8pt]{standalone}
\usepackage{ctex}
\usepackage{amsmath}
\usepackage{pgfplots}
\pgfplotsset{compat=1.18}
\usetikzlibrary{calc,arrows.meta,patterns}

\begin{document}
\begin{tikzpicture}[scale=1.0]

%% === 坐标系 ===
\draw[->, thick, gray!70] (-0.5, 0) -- (6.2, 0) node[right, font=\small] {对偶变量 $\lambda$};
\draw[->, thick, gray!70] (0, -0.5) -- (0, 5.2) node[above, font=\small] {目标值};

%% === 原始最优值 p*：水平实线 ===
\draw[red!70!black, very thick] (-0.3, 3.5) -- (6.0, 3.5);
\node[red!70!black, font=\small\bfseries, right] at (6.0, 3.5) {$p^*$（原始最优）};
\draw[red!70!black, dashed, thin] (0, 3.5) -- (0.1, 3.5);
\node[red!70!black, font=\small, left] at (0, 3.5) {$p^*$};

%% === 对偶函数 g(lambda)：抛物线形（凹函数）===
%% g(lambda) = -0.25*(lambda-4)^2 + 3.0（凹抛物线，最大值3.0在lambda=4）
\draw[blue!70!black, very thick, domain=0.2:7.8, samples=100]
  plot (\x, {-0.25*(\x-4)^2 + 3.0});
\node[blue!70!black, font=\small] at (1.0, 2.3) {$g(\lambda)$（对偶函数，凹）};

%% === 弱对偶区域（g <= p*）===
\fill[blue!10] (0.2, {-0.25*(0.2-4)^2+3.0})
  \foreach \x in {0.3,0.4,...,7.8} {-- (\x, {min(-0.25*(\x-4)^2+3.0, 3.5)})}
  -- (6.0, 3.5) -- (6.0, -0.5) -- (0.2, -0.5) -- cycle;

%% === 对偶最优 d*：g 的最大值处 ===
%% 最大值在 lambda*=4, d*=3.0
\coordinate (Dstar) at (4, 3.0);
\fill[blue!80!black] (Dstar) circle (3pt);
\node[blue!80!black, font=\small\bfseries, below right] at (Dstar) {$(\lambda^*, d^*)$};

%% 标注 d*
\draw[blue!70!black, dashed, thin] (4, 0) -- (4, 3.0);
\draw[blue!70!black, dashed, thin] (0, 3.0) -- (4, 3.0);
\node[blue!70!black, font=\small, left] at (0, 3.0) {$d^*$};
\node[blue!70!black, font=\small, below] at (4, 0) {$\lambda^*$};

%% === 对偶间隙（gap = p* - d*）===
\draw[<->, >=Stealth, very thick, green!50!black]
  (4.8, 3.0) -- (4.8, 3.5);
\node[green!50!black, font=\small\bfseries, right] at (4.9, 3.25)
  {对偶间隙 $p^*-d^*$};

%% === 弱对偶标注 ===
\node[font=\small, align=center, blue!50!black] at (1.5, 0.8)
  {弱对偶恒成立：\\$g(\lambda) \leq d^* \leq p^*$};

%% === 强对偶标注（Slater 条件满足时）===
\node[draw=red!40, fill=red!5, rounded corners, font=\small,
      align=center, inner sep=4pt] at (4.5, 4.8)
  {\textbf{强对偶}（Slater 条件满足）：\\
   $d^* = p^*$，对偶间隙 $= 0$};
\draw[->, red!50!black, dashed] (4.1, 4.5) -- (4.05, 3.6);

%% === 弱对偶注释 ===
\node[draw=blue!40, fill=blue!5, rounded corners, font=\small,
      align=center, inner sep=4pt] at (1.4, 4.3)
  {\textbf{弱对偶}（无条件成立）：\\$d^* \leq p^*$};

%% === 标题 ===
\node[font=\small\bfseries, align=center] at (3.0, -0.9)
  {对偶函数（凹）提供原始最优值的下界；Slater 条件满足时强对偶成立（$d^*=p^*$）};

\end{tikzpicture}
\end{document}
